%% S6 -- material moved from the main paper to meet its page budget. Nothing here is new: each item
%% comes from the main-paper section named in its lead-in, which keeps a pointer to it.
\section{Material moved from the main paper}
\label{app:moved}

This section carries reference material that the main paper summarizes in place and points to as
Supplement~S6. Each item comes from the section of the main paper named below, with its numbers
unchanged.

\subsection{Synthesis estimators}
\label{subsec:moved-estimators}

Table~\ref{tab:estimators} is the estimator table of \S4.7 of the main paper, which summarizes it in
one sentence.

\begin{table}[htbp]
  \footnotesize
  \caption{Synthesis estimators. Each question is answered by one named estimator, with the uncertainty or
  test that could overturn it. BH: Benjamini--Hochberg false-discovery control \citep{benjamini1995fdr},
  applied within each test family.}
  \label{tab:estimators}
  \begin{tabular}{@{}p{0.17\linewidth} p{0.44\linewidth} p{0.33\linewidth}@{}}
    \toprule
    Question & Estimator & Uncertainty and inference \\
    \midrule
    Field-level proportions & Stratified weighted proportion, 1,233 drawn systems & Design standard errors, the system as the sampling unit \\
    Coded-set proportions & Unweighted proportion, 1,256 systems & Wilson interval
      \citep{wilson1927interval} \\
    Diversity over time & Miller--Madow bias-corrected entropy in nats \citep{paninski2003entropy};
      year-by-dimension cells below 20 systems suppressed & Bootstrap interval; BH within the 34
      convergence verdicts \\
    Association & Bias-corrected Cram\'er's V \citep{cramer1946methods,bergsma2013cramersv} on complete
      pairs, no imputation (703 pairs estimable, 37 excluded); the matrix with \texttt{not\_reported} as a
      level reported beside it & BH across pairs \\
    Design families & Average-linkage clustering on a Gower distance \citep{gower1971similarity} over coded
      cells; 583 systems valued on at least half of the 37 dimensions, $k \leq 15$ & Silhouette
      \citep{rousseeuw1987silhouette}, bootstrap Jaccard stability \citep{hennig2007stability},
      selection-adjusted null \\
    Cross-paper comparison & Scores standardized within key (benchmark, split, base model), never pooled
      raw; the five contrasts of amendment~11 & Percentile bootstrap over keys; within-key permutation test,
      which governs on disagreement; minimum detectable effect at 80\% power; permutation design floor \\
    Within-study ablations & Relative change pooled per dimension by DerSimonian--Laird random effects
      \citep{dersimonian1986meta} & Confidence interval, $I^2$ \citep{higgins2002heterogeneity},
      prediction interval \citep{higgins2009predictioninterval}; bootstrap over comparison blocks \\
    Reporting bias & Exact binomial sign test on the contrasts; rank of the authors' own arm against the
      chance share implied by each block's arm count, computed exactly, per paper and one block per paper;
      trim-and-fill \citep{duval2000trimfill}; Copas-style null-fill \citep{copas2000funnel}; effect against
      benchmark size; Egger's test \citep{egger1997bias}, reported but not relied on because the variance
      is reconstructed from arm scores & Each pooled effect discounted by the larger of its trim-and-fill
      and one-null-per-contrast shrinkage \\
    \bottomrule
  \end{tabular}
\end{table}

\subsection{The full reporting checklist}
\label{subsec:moved-checklist}

Table~\ref{tab:checklist-full} is the complete reporting checklist of \S8.1 of the main paper, which
prints its five items with the highest weighted silence rate. It lists the twelve dimensions with the
highest weighted \texttt{not\_reported} rate, each turned into the minimum statement a paper or
repository should make.

\begin{table}[htbp]
\caption{A reporting checklist derived from the twelve dimensions with the highest weighted silence rate.
\emph{Silent} is the weighted census share of systems whose sources say nothing about the dimension
(design standard errors 1.5 to 3.1 percentage points). The rates measure documents, not systems.}
\label{tab:checklist-full}
\small
\begin{tabular}{llrp{0.50\linewidth}}
\toprule
Layer & Dimension & Silent & Minimum a paper or repository should state \\
\midrule
G & \texttt{network\_policy} & 93.3\% & Whether the agent can reach the network during a run, and if so what is reachable: open egress, an allowlist, or nothing. \\
H & \texttt{replayability} & 90.6\% & Whether a finished run can be replayed from stored artifacts, and whether fully, from inputs only, or not at all. \\
G & \texttt{filesystem\_access} & 87.5\% & The writable scope: the whole host, one working directory, a copy-on-write overlay, or read-only. \\
E & \texttt{rollback} & 86.4\% & What happens to partial edits when a step fails: snapshot restore, version-control revert, or nothing. \\
D & \texttt{state\_persistence} & 86.2\% & Whether a run can be resumed after interruption, and from what: full state, a checkpoint, or nothing. \\
B & \texttt{protocol\_standardization} & 84.0\% & Whether tools are exposed over a named standard protocol or over a bespoke interface. \\
F & \texttt{timeouts} & 82.4\% & The timeout enforced per tool call and per run, with the default value. \\
G & \texttt{permission\_model} & 81.4\% & Which actions require approval and from whom, or that none do. \\
F & \texttt{cost\_controls} & 80.2\% & The token, call, or currency budget enforced per run, and the behaviour when it is exhausted. \\
A & \texttt{context\_compaction} & 77.9\% & What happens when the context window fills: summarize, prune, checkpoint, truncate, or defer to the model. \\
B & \texttt{tool\_count} & 77.5\% & The number of tools registered in the default configuration. \\
C & \texttt{human\_in\_loop} & 76.7\% & Where a human can intervene, and whether the default configuration is autonomous. \\
\bottomrule
\end{tabular}
\end{table}

\subsection{Value distributions under weighting}
\label{subsec:moved-values}

Figure~\ref{fig:landscape_values} accompanies \S6.1 of the main paper, whose table of modal-value
changes it extends to every value of every coded dimension.

\begin{figure}[htbp]
\centering
\includegraphics[width=\linewidth]{descriptives_value_distributions}
\caption{Weighted value distributions for every coded dimension, with the unweighted share of the coded
set marked on each bar. The gap between bar and marker is the sampling distortion for that value, and it
is largest for \texttt{primary\_artifact}, where the modal value changes.}
\label{fig:landscape_values}
\end{figure}

\subsection{Silence per dimension}
\label{subsec:moved-underreporting}

Figure~\ref{fig:landscape_underreporting} accompanies \S6.2 of the main paper, whose table gives
the silence rate by layer and for the five least-documented dimensions; the figure gives it for
every dimension, unweighted against weighted.

\begin{figure}[htbp]
\centering
\includegraphics[width=0.9\linewidth]{descriptives_under_reporting_by_dimension}
\caption{Share of systems whose sources are silent, per dimension, unweighted against weighted. Weighting
raises the silence rate most for \texttt{open\_source}, \texttt{protocol\_standardization}, and
\texttt{execution\_isolation}, and lowers it for most control-loop dimensions. A review that codes only
visible systems understates the silence on what bounds an agent.}
\label{fig:landscape_underreporting}
\end{figure}

\subsection{Entropy trajectories}
\label{subsec:moved-trajectories}

Figure~\ref{fig:landscape_trajectories} accompanies \S6.3 of the main paper, whose table gives
the entropy change, interval and corrected $q$ for the ten dimensions whose interval excludes zero.

\begin{figure}[htbp]
\centering
\includegraphics[width=\linewidth]{descriptives_entropy_trajectories}
\caption{The largest entropy falls and rises, with the number of reporting systems at each point. The
$n$ labels are why falls are read more cautiously than rises: a fall with rising $n$ can reflect who
reports rather than what is built. Unweighted; the legend's verdicts are uncorrected, and all six plotted
dimensions survive the Benjamini--Hochberg correction.}
\label{fig:landscape_trajectories}
\end{figure}

\subsection{Where shared silence creates association}
\label{subsec:moved-assoc}

Figure~\ref{fig:landscape_assoc_delta} accompanies \S6.5 of the main paper. It locates, pair by pair,
the change in association when \texttt{not\_reported} is counted as a level rather than excluded.

\begin{figure}[htbp]
\centering
\includegraphics[width=\linewidth]{family_association_delta}
\caption{Corrected Cram\'er's V with \texttt{not\_reported} counted as a level minus the same coefficient
on complete pairs. Warm cells look more associated once silence is treated as a design choice; marked
cells change significance or move by at least 0.15. Most of the warmth is documentation habit, not design: 265 pairs
are significant only with silence as a level.}
\label{fig:landscape_assoc_delta}
\end{figure}
